% Generated draft: requires mathematical and source review.
\documentclass[11pt]{article}
\usepackage[margin=25mm]{geometry}
\usepackage{fontspec}
\setmainfont{Latin Modern Roman}
\usepackage{amsmath,amssymb,mathtools,unicode-math}
\setmathfont{Latin Modern Math}
\usepackage{enumitem,needspace,xurl,hyperref}
\setlength{\parindent}{0pt}
\setlength{\parskip}{5pt}
\setlength{\emergencystretch}{3em}
\newcommand{\N}{\mathbb{N}}
\newcommand{\Z}{\mathbb{Z}}
\newcommand{\Q}{\mathbb{Q}}
\newcommand{\R}{\mathbb{R}}
\newcommand{\C}{\mathbb{C}}
\newcommand{\F}{\mathbb{F}}
\newcommand{\A}{\mathbb{A}}
\newcommand{\PP}{\mathbb{P}}
\newcommand{\E}{\mathbb{E}}
\newcommand{\eps}{\varepsilon}
\newcommand{\dd}{\,\mathrm{d}}
\begin{document}
\section*{UnsolvedMath ID: 20000331}
\textbf{Source problem code:} catalogue entry
\subsection{Definitions and mathematical statement}
Viehweg's source question asks whether the number of families of minimal manifolds with Hilbert polynomial \(H\), over a curve of genus \(q\) with \(s\) singular fibers, has a universal bound depending only on \(H,q,s\), not further fiber geometry. The record does not define the equivalence relation on families, whether the pointed base is fixed, whether non-isotriviality is required, or which polarization determines \(H\). Counting individual isomorphism classes and counting deformation types give different questions.

The attached report proposes a counterexample only to the literal isomorphism-class reading, even with a fixed pointed base and non-isotriviality. Choose one generic plane-quartic Lefschetz pencil \(\mathcal C\to B=\PP^1\), with its fixed singular set \(S\subset B\) of exactly \(27\) points. Let \(D_j\) range over infinitely many genus-four curves whose Jacobians are simple and pairwise nonisogenous. The proposed surface families are
\[
\mathcal X_j=\mathcal C\times D_j\longrightarrow B.
\]
Their smooth fibers are products of a genus-three curve and \(D_j\), hence are canonically polarized surfaces. The report asserts that all these families are non-isotrivial, have exactly the same bad set \(S\), and have the same canonical Hilbert polynomial
\[
H(m)=\chi(X_{j,b},\omega_{X_{j,b}}^{\otimes m})=6(2m-1)^2.
\]
It claims they are pairwise nonisomorphic over the fixed base, using Albanese/Jacobian isogeny cancellation, while a connected level-moduli family shows they all have a single deformation type. Thus the counterexample fixes \(q=0\), \(s=27\), the whole pointed curve \((B,S)\), non-isotriviality, and \(H\), not merely their numerical bounds. No explicit pencil or infinite curve sequence is given in the supplied summary.

The report says this is compatible with a Kov\'acs--Lieblich uniform theorem for deformation types, which it calls the accepted canonically polarized interpretation of the AIM question. The full theorem scope and the broader minimal-manifold setting are not supplied, so neither a new contradiction to that theorem nor a complete uniformity theorem for all minimal manifolds is asserted here.

The extracted record also contains adjacent Section 3 GIT background, not an additional formulation of Viehweg's question. It says GIT divides Hilbert-scheme points into stable, semistable, and unstable strata; the stable quotient is geometric, while the semistable good quotient can identify nonisomorphic varieties whose orbit closures meet. It emphasizes that GIT stability differs from MMP stability: some MMP-stable singularities are asymptotically GIT-unstable, and some GIT-semistable points are not MMP-stable. These contextual claims do not change the required equivalence relation for counting families.
\subsection{Short English statement}
Ask for a bound on families depending only on a Hilbert polynomial and base-curve data. The report gives infinitely many isomorphism classes with one fixed base and bad set, but puts them in one deformation type, preserving the distinction needed for the cited uniform theorem.
\subsection{Sources}
\begin{itemize}
\item AIM workshop ``Compact moduli spaces and birational geometry,'' Section 2, ``Moduli spaces of manifolds,'' Viehweg's question numbered 4 in the extract: \url{https://aimath.org/WWN/birational/birational.pdf}. No PDF was accessed here.
\item Local canonical registry \texttt{aim-algebraic-geometry-notes.json}, \texttt{notes[330]}, tag \texttt{problem}, source item recorded as 2; permanent UnsolvedMath ID \(20000331\), source code \texttt{catalogue entry}. The extract includes a running header naming Michael A. van Opstall and neighboring GIT background. Original remarks and literature field are empty.
\item Supplied research attempt 1 attributes deformation-type uniformity to Kov\'acs--Lieblich and proposes the fixed-base product counterexample. No full authoritative bibliographic references or explicit pencil/curve data are supplied. The report is machine generated, not peer reviewed, and requires independent verification before citation.
\end{itemize}
\end{document}
